% ═══════════════════════════════════════════════════════════════════════════
% KLASSENARBEIT: Examen Unidad 1 + 2 (Klasse 9) - Version 3.0
% ═══════════════════════════════════════════════════════════════════════════

\documentclass[a4paper,11pt]{article}

\usepackage[utf8]{inputenc}
\usepackage[T1]{fontenc}
\usepackage[ngerman]{babel}
\usepackage{helvet}
\renewcommand{\familydefault}{\sfdefault}

\usepackage[margin=18mm]{geometry}
\usepackage{xcolor}
\usepackage{tikz}
\usepackage{enumitem}
\usepackage{array}
\usepackage{booktabs}
\usepackage{multicol}
\usepackage{tabularx}
\usepackage{tcolorbox}
\usepackage{amssymb}
\usepackage{fancyhdr}
\usepackage{lastpage}

% Fachfarbe Spanisch
\definecolor{spanisch}{HTML}{c41e3a}
\definecolor{hellgrau}{HTML}{f5f5f5}
\definecolor{mittelgrau}{HTML}{888888}
\definecolor{wahlfarbe}{HTML}{2c5aa0}

% Lücke
\newcommand{\luecke}[1]{\underline{\hspace{#1}}}

% Punktebox (rechts)
\newcommand{\punkte}[1]{\hfill\textcolor{mittelgrau}{\small /\,#1}}

% Wahlaufgaben-Box
\newcommand{\wahlhinweis}{\noindent\fcolorbox{wahlfarbe}{wahlfarbe!10}{\parbox{0.96\textwidth}{\small\color{wahlfarbe}\textbf{Elige una opción} -- Wähle eine Variante (A oder B).}}\vspace{2mm}}

% Aufgaben-Umgebung
\newcounter{aufgabe}
\newenvironment{aufgabe}[2]{%
    \refstepcounter{aufgabe}
    \vspace{4mm}
    \noindent\textbf{\color{spanisch}Ejercicio \theaufgabe: #1 \punkte{#2}}
    \vspace{1.5mm}
    \par
}{%
    \vspace{2mm}
}

\setlength{\parindent}{0pt}
\setlength{\parskip}{0pt}

% Seitennummerierung
\pagestyle{fancy}
\fancyhf{}
\renewcommand{\headrulewidth}{0pt}
\fancyfoot[C]{\small\color{mittelgrau}Seite \thepage\ von \pageref{LastPage}}

\begin{document}

% ═══════════════════════════════════════════════════════════════════════════
% KOPFBEREICH
% ═══════════════════════════════════════════════════════════════════════════

\noindent
\begin{tabularx}{\textwidth}{Xr}
    {\Large\bfseries\color{spanisch}Examen: Unidad 1 + 2} & {\small Klasse 9 -- Spanisch} \\
\end{tabularx}

\vspace{1mm}
\noindent\rule{\textwidth}{0.4pt}
\vspace{2mm}

\noindent
{\small Nombre: \luecke{6cm} \hfill Fecha: \luecke{3cm} \hfill Punkte: \luecke{1.5cm} / 49}

\vspace{3mm}

% Notentabelle 14-NP (Sek I, Klasse 9)
\noindent
{\footnotesize
\begin{tabular}{|c|c|c|c|c|c|c|c|c|c|c|c|c|c|c|}
\hline
\textbf{NP} & 14 & 13 & 12 & 11 & 10 & 9 & 8 & 7 & 6 & 5 & 4 & 3 & 2 & 1 \\
\hline
\textbf{ab} & 49 & 47,5 & 44,5 & 42,5 & 39,5 & 36 & 33 & 29,5 & 26,5 & 23 & 20 & 16,5 & 13,5 & 10 \\
\hline
\end{tabular}
}

\vspace{4mm}

% ═══════════════════════════════════════════════════════════════════════════
% AUFGABE 1: Vocabulario im Kontext (AFB II)
% ═══════════════════════════════════════════════════════════════════════════

\begin{aufgabe}{Vocabulario en contexto}{6}

Lee las situaciones y completa con la palabra española correcta (con artículo).

{\small\textit{Lies die Situationen und ergänze das passende spanische Wort mit Artikel.}}

\vspace{2mm}

\begin{enumerate}[itemsep=4mm]
    \item Du möchtest im Restaurant bezahlen. Du fragst den Kellner nach \textit{la/el} \luecke{3cm}.
    \item Der Kellner bringt dir Suppe. Du brauchst zum Essen \textit{la/el} \luecke{3cm}.
    \item Du ziehst in eine neue Wohnung. Die Milch stellst du in \textit{la/el} \luecke{3cm}.
    \item Im Wohnzimmer hängst du \textit{la/el} \luecke{3cm} an die Wand.
    \item Du möchtest dich nach dem Duschen sehen. Du schaust in \textit{la/el} \luecke{3cm}.
    \item Deine Bücher ordnest du in \textit{la/el} \luecke{3cm}.
\end{enumerate}

\end{aufgabe}

% ═══════════════════════════════════════════════════════════════════════════
% AUFGABE 2: Verbos pedir/servir - Fehlerkorrektur (AFB III)
% ═══════════════════════════════════════════════════════════════════════════

\begin{aufgabe}{Corrige los errores -- pedir y servir}{5}

Encuentra el error en cada frase y escribe la forma correcta del verbo.

{\small\textit{Finde den Fehler und schreibe die korrigierte Verbform.}}

\vspace{2mm}

\begin{enumerate}[itemsep=4mm]
    \item En el restaurante, yo siempre \textbf{*pedo} una paella.

    Corrección / Korrektur: \luecke{3cm}

    \item Los camareros \textbf{*serven} muy rápido.

    Corrección / Korrektur: \luecke{3cm}

    \item ¿Qué \textbf{*pedís} tú de postre?

    Corrección / Korrektur: \luecke{3cm}

    \item Nosotros \textbf{*servemos} las bebidas primero.

    Corrección / Korrektur: \luecke{3cm}

    \item Mi hermana nunca \textbf{*pede} carne.

    Corrección / Korrektur: \luecke{3cm}
\end{enumerate}

\end{aufgabe}

\newpage
% ═══════════════════════════════════════════════════════════════════════════
% AUFGABE 3: Conocer + a personal (AFB II)
% ═══════════════════════════════════════════════════════════════════════════

\begin{aufgabe}{El verbo conocer -- ¿con o sin \textit{a}?}{4}

Completa con la forma correcta de \textbf{conocer} y decide: ¿\textbf{a} o \textbf{--}?

{\small\textit{Ergänze die richtige Form von conocer und entscheide: a oder -- (keine Präposition)?}}

\vspace{2mm}

\begin{enumerate}[itemsep=4mm]
    \item Yo \luecke{2.5cm} \luecke{1cm} tu hermano.
    \item ¿Tú \luecke{2.5cm} \luecke{1cm} Barcelona bien?
    \item Nosotros \luecke{2.5cm} \luecke{1cm} la profesora nueva.
    \item Ellos \luecke{2.5cm} \luecke{1cm} muchos restaurantes.
\end{enumerate}

\end{aufgabe}

% ═══════════════════════════════════════════════════════════════════════════
% AUFGABE 4: Pretérito perfecto - gemischt (AFB I/II)
% ═══════════════════════════════════════════════════════════════════════════

\begin{aufgabe}{El pretérito perfecto}{6}

Forma el pretérito perfecto con los verbos entre paréntesis.

{\small\textit{Bilde das pretérito perfecto mit den Verben in Klammern.}}

\vspace{2mm}

\begin{enumerate}[itemsep=4mm]
    \item Esta mañana yo \luecke{4cm} mucho. (trabajar)
    \item ¿Qué \luecke{4cm} (tú) hoy en clase? (hacer)
    \item Mis padres nunca \luecke{4cm} a Alemania. (viajar)
    \item Yo todavía no \luecke{4cm} la película. (ver)
    \item Nosotros \luecke{4cm} una carta a la abuela. (escribir)
    \item Hoy María no \luecke{4cm} nada. (comer)
\end{enumerate}

\end{aufgabe}


% ═══════════════════════════════════════════════════════════════════════════
% AUFGABE 5: Direktobjektpronomen im Dialog (AFB II)
% ═══════════════════════════════════════════════════════════════════════════

\begin{aufgabe}{Los pronombres de objeto directo}{5}

Completa el diálogo con los pronombres correctos (lo/la/los/las).

{\small\textit{Ergänze den Dialog mit den passenden Direktobjektpronomen.}}

\vspace{2mm}

\begin{tcolorbox}[colback=hellgrau, colframe=mittelgrau, boxrule=0.5pt, arc=0pt]
\textbf{Mamá:} ¿Dónde está la lámpara nueva?\\[1mm]
\textbf{Pablo:} \luecke{1.2cm} he puesto en el dormitorio.\\[2mm]
\textbf{Mamá:} ¿Y los libros de la escuela?\\[1mm]
\textbf{Pablo:} \luecke{1.2cm} he dejado en la estantería.\\[2mm]
\textbf{Mamá:} ¿Has visto mis gafas?\\[1mm]
\textbf{Pablo:} Sí, \luecke{1.2cm} he visto en la cocina.\\[2mm]
\textbf{Mamá:} ¿Y el sofá nuevo? ¿Ya está en el salón?\\[1mm]
\textbf{Pablo:} Sí, papá y yo \luecke{1.2cm} hemos subido esta mañana.\\[2mm]
\textbf{Mamá:} ¡Perfecto! ¿Y las plantas del balcón?\\[1mm]
\textbf{Pablo:} \luecke{1.2cm} he regado antes de salir.
\end{tcolorbox}

\end{aufgabe}
\newpage
% ═══════════════════════════════════════════════════════════════════════════
% AUFGABE 6: Bewegungsverben + Präpositionen (AFB II)
% ═══════════════════════════════════════════════════════════════════════════

\begin{aufgabe}{Los verbos de movimiento -- elige la preposición}{5}

Elige la preposición correcta.

{\small\textit{Wähle die richtige Präposition.}}

\vspace{2mm}

\begin{tabular}{|p{9cm}|c|c|c|c|}
\hline
\textbf{Frase / Satz} & \textbf{de} & \textbf{a} & \textbf{en} & \textbf{del/al} \\
\hline
Bajo las cajas \rule{0.8cm}{0.4pt} sótano. & $\square$ & $\square$ & $\square$ & $\square$ \\[2mm]
\hline
Llevo el televisor \rule{0.8cm}{0.4pt} dormitorio. & $\square$ & $\square$ & $\square$ & $\square$ \\[2mm]
\hline
Saco la leche \rule{0.8cm}{0.4pt} la nevera. & $\square$ & $\square$ & $\square$ & $\square$ \\[2mm]
\hline
Pongo los libros \rule{0.8cm}{0.4pt} la estantería. & $\square$ & $\square$ & $\square$ & $\square$ \\[2mm]
\hline
Subo las maletas \rule{0.8cm}{0.4pt} piso. & $\square$ & $\square$ & $\square$ & $\square$ \\[2mm]
\hline
\end{tabular}

\end{aufgabe}



% ═══════════════════════════════════════════════════════════════════════════
% AUFGABE 7: Leseverstehen (AFB II/III)
% ═══════════════════════════════════════════════════════════════════════════

\begin{aufgabe}{Comprensión de lectura}{8}

Lee el texto y contesta las preguntas \textbf{en español}.

{\small\textit{Lies den Text und beantworte die Fragen auf Spanisch.}}

\vspace{2mm}

\begin{tcolorbox}[colback=hellgrau, colframe=mittelgrau, boxrule=0.5pt, arc=0pt, title={\small\bfseries Una cena especial}]
{\small
Este fin de semana mi familia y yo hemos ido a un restaurante italiano para celebrar el cumpleaños de mi madre. El restaurante se llama ``La Dolce Vita'' y está en el centro de la ciudad.

Cuando hemos llegado, el camarero nos ha llevado a una mesa cerca de la ventana. Mi padre ha pedido una pizza con jamón, mi madre ha pedido pasta con mariscos y yo he pedido una ensalada grande porque no he tenido mucha hambre.

El camarero ha servido la comida muy rápido -- solo hemos esperado quince minutos. Todo ha estado delicioso. De postre, hemos pedido un tiramisú para compartir. Al final, mi padre ha pedido la cuenta. Ha costado cuarenta y cinco euros, y mi padre ha dejado cinco euros de propina porque el servicio ha sido muy bueno.

Mi madre ha dicho: ``¡Ha sido la mejor cena de cumpleaños!'' Yo pienso que tiene razón -- el restaurante ha sido perfecto.
}
\end{tcolorbox}

\vspace{3mm}

\begin{enumerate}[itemsep=5mm]
    \item ¿Por qué han ido al restaurante? (1P)

    \luecke{12cm}

    \item ¿Qué ha pedido cada persona de la familia? Escribe los tres platos. (3P)

    \luecke{12cm}

    \item ¿Por qué ha pedido el/la narrador/a una ensalada? (1P)

    \luecke{12cm}

    \item ¿Cuánto ha costado la cena en total (con propina)? Calcula. (1P)

    \luecke{12cm}

    \item ¿Crees que la familia va a volver a este restaurante? ¿Por qué (no)? (2P)

    \luecke{12cm}

\end{enumerate}

\end{aufgabe}

\newpage

% ═══════════════════════════════════════════════════════════════════════════
% AUFGABE 8: Schreibaufgabe - Wahlaufgabe (AFB III)
% ═══════════════════════════════════════════════════════════════════════════

\begin{aufgabe}{Producción escrita}{10}

\wahlhinweis

\textbf{Variante A: Diálogo en el restaurante}

Escribe un diálogo entre un/a camarero/a y un/a cliente.

{\small\textit{Schreibe einen Dialog zwischen Kellner/in und Gast.}}

\vspace{1mm}

El diálogo debe incluir / Der Dialog muss enthalten:
\begin{itemize}[itemsep=0mm, topsep=1mm]
    \item Primer plato, segundo plato y bebida
    \item Una pregunta sobre el postre
    \item Pedir la cuenta
    \item Mínimo 10 intercambios (5 del camarero, 5 del cliente)
\end{itemize}

\vspace{2mm}

\begin{tabular}{p{2.5cm}p{12cm}}
\textbf{Camarero/a:} & \rule{12cm}{0.4pt} \\[5mm]
\textbf{Cliente:} & \rule{12cm}{0.4pt} \\[5mm]
\textbf{Camarero/a:} & \rule{12cm}{0.4pt} \\[5mm]
\textbf{Cliente:} & \rule{12cm}{0.4pt} \\[5mm]
\textbf{Camarero/a:} & \rule{12cm}{0.4pt} \\[5mm]
\textbf{Cliente:} & \rule{12cm}{0.4pt} \\[5mm]
\textbf{Camarero/a:} & \rule{12cm}{0.4pt} \\[5mm]
\textbf{Cliente:} & \rule{12cm}{0.4pt} \\[5mm]
\textbf{Camarero/a:} & \rule{12cm}{0.4pt} \\[5mm]
\textbf{Cliente:} & \rule{12cm}{0.4pt} \\
\end{tabular}

\newpage

\textbf{Variante B: Mi nueva casa}

Describe tu casa o piso (real o imaginario).

{\small\textit{Beschreibe dein Haus oder deine Wohnung (real oder erfunden).}}

\vspace{1mm}

El texto debe incluir / Der Text muss enthalten:
\begin{itemize}[itemsep=0mm, topsep=1mm]
    \item Mínimo 4 habitaciones con sus muebles
    \item Mínimo 80 palabras
\end{itemize}

\vspace{3mm}

\noindent\rule{\textwidth}{0.4pt}\\[6mm]
\rule{\textwidth}{0.4pt}\\[6mm]
\rule{\textwidth}{0.4pt}\\[6mm]
\rule{\textwidth}{0.4pt}\\[6mm]
\rule{\textwidth}{0.4pt}\\[6mm]
\rule{\textwidth}{0.4pt}\\[6mm]
\rule{\textwidth}{0.4pt}\\[6mm]
\rule{\textwidth}{0.4pt}\\[6mm]
\rule{\textwidth}{0.4pt}\\[6mm]
\rule{\textwidth}{0.4pt}\\[6mm]
\rule{\textwidth}{0.4pt}\\[6mm]
\rule{\textwidth}{0.4pt}\\[6mm]
\rule{\textwidth}{0.4pt}\\[6mm]
\rule{\textwidth}{0.4pt}\\[6mm]
\rule{\textwidth}{0.4pt}\\[6mm]
\rule{\textwidth}{0.4pt}

\end{aufgabe}

\vfill

\noindent\rule{\textwidth}{0.4pt}

\noindent{\small\color{mittelgrau}Bewertung Ejercicio 8: Inhalt (3P) | Grammatik (3P) | Kohäsion/Kohärenz (2P) | Rechtschreibung (2P)}

\end{document}
